% !PRINTED 558
\weberpage{558}

If we sum these formulae once again with respect to $y$ from $1$ to
$\infty$, and make use of the formulae which are partly known and
partly follow from \S\,24, (11):

\[
\sum^{y}\frac{1}{y^2}=\frac{\pi^2}{6},\quad
\sum^{y}\frac{(-1)^y}{y^2}=-\frac{\pi^2}{12},\quad
\sum\frac{(-1)^y}{y}=-\log 2,\,.
\]

\[
\sum^{n}\sum^{y}\frac{q^{(2n-1)y}}{y}
=-\log\Pi(1-q^{2n-1})=-\log f_1(\omega)-\frac{\pi i\omega}{24},
\]

\[
\sum^{n}\sum^{y}\frac{(-1)^y q^{2ny}}{y}
=-\log\Pi(1+q^{2n})=-\log f_2(\omega)+\frac{\pi i\omega}{12}-\tfrac{1}{2}\log 2,
\]

\[
\sum^{n}\sum^{y}\frac{(-1)^y q^{(2n-1)y}}{y}
=-\log\Pi(1+q^{2n-1})=-\log f(\omega)-\frac{\pi i\omega}{24},
\]

then there results:

\weq{5}{\begin{aligned}
\sum_{1,\infty}^{y}\sum_{1,\infty}^{x}\frac{(-1)^{x+y}\omega}{\omega^2y^2-x^2}
&=\frac{\pi^2}{24\omega}-\frac{\pi^2\omega}{24}+\pi i\log f(\omega),\\[4pt]
\sum^{y}\sum^{x}\frac{(-1)^x\omega}{\omega^2y^2-x^2}
&=-\frac{\pi^2}{12\omega}-\frac{\pi^2\omega}{24}+\pi i\log f_1(\omega),\\[4pt]
\sum^{y}\sum^{x}\frac{(-1)^y\omega}{\omega^2y^2-x^2}
&=\frac{\pi^2}{24\omega}+\frac{\pi^2\omega}{12}+\pi i\log f_2(\omega)\,.
\end{aligned}}

The double sums on the left side of these formulae are to be understood
in such a way that the summation with respect to $x$ is to be carried
out first, and then the summation with respect to $y$. But we may also
sum in such a way that we take

\weq{6}{
0<x<m,\quad 0<y<n,\quad n=\infty\,,\quad \frac{m}{n}=\infty
}

and then let $m$ and $n$ grow to infinity so that $m:n$ becomes
infinite. That both give the same can be shown in various ways, for
example by comparing the sums with integrals.

If, then, we set, as in \S\,147, (10), binding $x$, $y$ to the
condition (6):

\[
\begin{aligned}
\sum^{x,y}\frac{\omega}{\omega^2y^2-x^2}&=S_0 & x,\,y&\ \text{even,}\\
&=S_1 & x,\,y&\ \text{odd,}\\
&=S_2 & x\ \text{odd,}\ y&\ \text{even,}\\
&=S_3 & x\ \text{even,}\ y&\ \text{odd,}
\end{aligned}
\]

it follows from (5):
%CONT>>
% [sic] The printed147,(10) reference at the S0,S1,S2,S3 definition names preceding log formulae instead of definition(11). Retained.
